\documentclass[11pt]{article}
\usepackage[margin=1in]{geometry}
\usepackage{amsmath,amssymb,amsthm}
\usepackage{graphicx}
\usepackage[hidelinks]{hyperref}

\newtheorem{theorem}{Theorem}
\newtheorem{remark}{Remark}
\newcommand{\Hh}{\mathcal H}
\newcommand{\N}{\mathbb N}
\newcommand{\ess}{\mathrm{ess}}
\newcommand{\diag}{\operatorname{diag}}
\newcommand{\spn}{\operatorname{span}}
\newcommand{\dist}{\operatorname{dist}}

\title{A Fixed-Section Counterexample for the Infinite-Dimensional QR Algorithm}
\author{Agent lane 5, run \texttt{fa\_banach\_001}}
\date{2026-07-19}

\begin{document}
\maketitle

\section*{Source and Target}

The source paper is M. J. Colbrook and A. C. Hansen,
\emph{On the infinite-dimensional QR algorithm}, arXiv:2011.08172
\cite{colbrook_hansen_iqr}.  In the concluding remarks, the authors write that
their examples lead them to conjecture that, for normal operators, the IQR
algorithm recovers the extreme parts of the essential spectrum when the set of
extremal points of the essential spectrum has size one.  In the source TeX this
is at lines 2167--2171; in the locally compiled source PDF it is on page 42.

\begin{center}
\includegraphics[width=0.95\linewidth]{figures/open_problem_crop.png}
\end{center}

The counterexample below addresses the natural fixed-finite-section reading of
that sentence: after choosing a sufficiently large finite \(m\), the limit
\(n\to\infty\) of the spectra of
\[
  P_m Q_n^* T Q_n\vert_{P_m\Hh}
\]
should recover the unique extremal point of the essential spectrum.  The example
shows that no finite \(m\), however large, can do this in general.

\section*{Proof Intuition}

The IQR iterates do not create spectral support in the first \(m\) columns.  The
source paper records the identity
\[
  \spn\{Q_n e_j:1\leq j\leq m\}
  =
  \spn\{T^n e_j:1\leq j\leq m\}.
\]
If the first \(m\) basis vectors sit inside a finite reducing subspace \(F\) of
\(T\), then the IQR \(m\)-section is only a compression of \(T\) to a subspace
of \(F\).  Thus it can never see spectral points that occur only as an
asymptotic tail of later reducing blocks.  Taking non-diagonal \(2\times2\)
blocks whose eigenvalues increase to \(1\) gives \(\sigma_{\ess}(T)=\{1\}\), but
each fixed finite leading section remains separated from \(1\).

\section*{Counterexample}

\begin{theorem}
There exists an invertible normal operator \(T\in\mathcal B(\ell^2(\N))\) such
that:
\begin{enumerate}
\item \(T\) is self-adjoint, positive, and non-diagonal in the canonical basis;
\item \(\sigma_{\ess}(T)=\{1\}\), so the essential spectrum has exactly one
point of maximal modulus;
\item \(T\) has no eigenvalue of modulus greater than \(r_{\ess}(T)=1\);
\item for every \(m\in\N\) there is a number \(\varepsilon_m>0\) such that, for
every IQR \(Q\)-sequence \((Q_n)\) of \(T\) and every \(n\in\N\),
\[
  \dist\!\left(1,\sigma(P_mQ_n^*TQ_n\vert_{P_m\ell^2})\right)
  \geq \varepsilon_m.
\]
\end{enumerate}
Consequently the unique extremal point \(1\in\sigma_{\ess}(T)\) is not recovered
by the fixed-\(m\), \(n\to\infty\) IQR sections for any sufficiently large finite
\(m\).
\end{theorem}

\begin{proof}
For \(k\geq1\), set
\[
  a_k=1-\frac1{2k},\qquad b_k=1-\frac1{2k+1}.
\]
Then
\[
  0<a_1<b_1<a_2<b_2<\cdots<1,\qquad a_k,b_k\to1.
\]
Let
\[
  U=\frac1{\sqrt2}
  \begin{pmatrix}
  1&1\\
  1&-1
  \end{pmatrix}
\]
and define
\[
  B_k=U^*\diag(a_k,b_k)U
  =
  \begin{pmatrix}
  (a_k+b_k)/2&(a_k-b_k)/2\\
  (a_k-b_k)/2&(a_k+b_k)/2
  \end{pmatrix}.
\]
Each \(B_k\) is a non-diagonal real symmetric positive matrix with eigenvalues
\(a_k\) and \(b_k\).  With respect to the canonical decomposition
\[
  \ell^2(\N)=\bigoplus_{k=1}^{\infty}
  \spn\{e_{2k-1},e_{2k}\},
\]
define
\[
  T=\bigoplus_{k=1}^{\infty} B_k.
\]
Then \(T\) is bounded, self-adjoint, positive, and invertible, since
\(\inf_k a_k=a_1=1/2\).  It is not diagonal in the canonical basis because
\((a_k-b_k)/2\neq0\) for every \(k\).

Moreover,
\[
  T-I=\bigoplus_{k=1}^{\infty}(B_k-I_2)
\]
is compact, because \(\|B_k-I_2\|=1-a_k\to0\).  In the Calkin algebra, the image
of \(T\) is therefore the image of \(I\), and hence
\[
  \sigma_{\ess}(T)=\{1\}.
\]
The point \(1\) is the unique point of the essential spectrum and the unique
point of maximal modulus in it.  The eigenvalues of \(T\) are exactly the
numbers \(a_k,b_k\), all of which are strictly smaller than \(1\).  Thus there
are no eigenvalues with modulus greater than \(r_{\ess}(T)=1\).

Fix \(m\in\N\) and let \(K=\lceil m/2\rceil\).  Put
\[
  F_K=\bigoplus_{k=1}^{K}\spn\{e_{2k-1},e_{2k}\}.
\]
This is a finite-dimensional reducing subspace for \(T\), and
\(\spn\{e_1,\ldots,e_m\}\subset F_K\).  Therefore
\[
  T^n\spn\{e_1,\ldots,e_m\}\subset F_K
  \qquad(n\geq0).
\]
For the IQR \(Q\)-sequence, the source paper's identity
\[
  \spn\{Q_n e_j:1\leq j\leq m\}
  =
  \spn\{T^n e_j:1\leq j\leq m\}
\]
follows from \(T^n=Q_nR_n\) and the fact that \(R_n\) is upper triangular with
non-zero diagonal.  Hence
\[
  V_{m,n}:=\spn\{Q_ne_1,\ldots,Q_ne_m\}\subset F_K
  \qquad(n\geq0).
\]

The matrix \(P_mQ_n^*TQ_n\vert_{P_m\ell^2}\) is the matrix of the compression of
the self-adjoint operator \(T\) to the \(m\)-dimensional subspace \(V_{m,n}\),
written in the orthonormal basis \(Q_ne_1,\ldots,Q_ne_m\).  Since
\(V_{m,n}\subset F_K\), the min-max principle gives
\[
  \sigma(P_mQ_n^*TQ_n\vert_{P_m\ell^2})
  \subset [a_1,b_K].
\]
Consequently
\[
  \dist\!\left(1,\sigma(P_mQ_n^*TQ_n\vert_{P_m\ell^2})\right)
  \geq 1-b_K
  = \frac1{2K+1}
  >0
\]
for every \(n\).  Taking \(\varepsilon_m=1/(2\lceil m/2\rceil+1)\) proves the
claim.
\end{proof}

\section*{Interpretive Scope}

This packet does not rule out weaker double-limit statements in which \(m\) is
allowed to tend to infinity after, before, or together with \(n\).  Indeed, in
the example above, \(b_{\lceil m/2\rceil}\to1\) as \(m\to\infty\).  What fails is
the fixed finite-section conclusion suggested by the source phrase
``in the limit \(n\to\infty\) for large enough \(m\)'': no finite \(m\) has
\(n\)-iterates whose \(m\)-section can reach the unique essential extremal point.

\section*{Verification Notes}

The proof is analytic and uses no numerical hypothesis.  The accompanying
script \path{code/verify_block_counterexample.py} checks finite block
eigenvalues, the non-diagonal form of the blocks, and the predicted positive gaps
\[
  1-b_{\lceil m/2\rceil}.
\]
It is only a sanity check; the proof above is independent of floating-point
computation.

\section*{Novelty Check}

Before promotion, the run indexes were searched for arXiv:2011.08172,
\emph{On the infinite-dimensional QR algorithm}, IQR, QR algorithm, QR
iteration, and Hessenberg; no duplicate solution packet or attempt for this
paper was found.  A bounded web search on 2026-07-19 for exact title/id phrases
and the phrases ``IQR algorithm'', ``normal operators'', ``essential spectrum'',
and ``extremal points'' found the source paper but no later answer or
counterexample to this exact conjectural fixed-section statement.  The novelty
confidence is therefore moderate: the construction is elementary, but no
pre-existing answer was found in the checked sources.

\section*{References}

\begin{thebibliography}{1}

\bibitem{colbrook_hansen_iqr}
M. J. Colbrook and A. C. Hansen,
\emph{On the infinite-dimensional QR algorithm},
Numerische Mathematik 143 (2019), 17--83; arXiv:2011.08172.

\end{thebibliography}

\end{document}
